% =================================================
% Ученический рабочий лист
% =================================================
\worksheettitle
  {Рациональное уравнение с модулем и параметром}
  {Рабочий лист: открываем решение шаг за шагом}
\studentnameblock

\begin{problembox}
  Найдите все значения параметра \(a\), при каждом из которых уравнение
  \[
    \frac{|3x|-2x-2-a}{x^2-2x-a}=0
  \]
  имеет ровно два различных корня.
\end{problembox}

\begin{thinkbox}[Цель работы]
  Не получить готовую формулу, а построить надёжный маршрут решения:
  найти кандидатов из числителя, проверить условия ветвей модуля и убедиться,
  что каждый кандидат входит в область определения дроби.
\end{thinkbox}

\section{Что означает «дробь равна нулю»}

Сравните два уравнения:
\[
  \frac{x-2}{x+1}=0,
  \qquad
  \frac{x-2}{x-2}=0.
\]

\begin{fillbox}[Проверьте значение \(x=2\)]
  \begin{center}
    \renewcommand{\arraystretch}{1.42}
    \begin{tabularx}{\textwidth}{
      >{\centering\arraybackslash}p{0.22\textwidth}
      >{\centering\arraybackslash}p{0.22\textwidth}
      >{\centering\arraybackslash}p{0.22\textwidth}
      >{\centering\arraybackslash}X}
      \toprule
      \rowcolor{TableHeader}
      Выражение & Числитель & Знаменатель & Является ли \(x=2\) корнем? \\
      \midrule
      \(\dfrac{x-2}{x+1}\) & \answerblank[18mm] & \answerblank[18mm] & \answerblank[25mm] \\
      \addlinespace[0.25em]
      \(\dfrac{x-2}{x-2}\) & \answerblank[18mm] & \answerblank[18mm] & \answerblank[25mm] \\
      \bottomrule
    \end{tabularx}
  \end{center}
\end{fillbox}

Сформулируйте правило:
\[
  \frac{A(x)}{B(x)}=0
  \quad\Longleftrightarrow\quad
  \begin{cases}
    A(x)=\answerblank[18mm],\\[1mm]
    B(x)\answerblank[10mm]0.
  \end{cases}
\]

\begin{checkpointbox}[Два слоя основной задачи]
  Дополните схему словами «два различных нуля» и
  «не обращают знаменатель в ноль».
  \parameterlogicworksheetfigure
\end{checkpointbox}

\begin{thinkbox}[Прогноз до вычислений]
  После раскрытия модуля числитель задаётся двумя
  \answerblank[28mm] выражениями. Поэтому на каждой ветви может появиться
  не более \answerblank[10mm] корня, а всего — не более
  \answerblank[10mm] корней.
\end{thinkbox}

\section{Исследуем нули числителя}

Уравнение числителя перепишем так:
\[
  3|x|-2x=a+2.
\]

\begin{fillbox}[Раскройте модуль по двум ветвям]
  \begin{enumerate}
    \item При \(x\ge0\): \(|3x|=\answerblank[18mm]\), поэтому
      \[
        3x-2x=a+2
        \quad\Longrightarrow\quad
        x_1=\answerblank[25mm].
      \]
      Условие выбранной ветви:
      \[
        a+2\ge0
        \quad\Longleftrightarrow\quad
        a\ge\answerblank[14mm].
      \]

    \item При \(x<0\): \(|3x|=\answerblank[18mm]\), поэтому
      \[
        -3x-2x=a+2
        \quad\Longrightarrow\quad
        x_2=\answerblank[30mm].
      \]
      Условие выбранной ветви:
      \[
        -\frac{a+2}{5}<0
        \quad\Longleftrightarrow\quad
        a>\answerblank[14mm].
      \]
  \end{enumerate}
\end{fillbox}

Соберите результаты в таблицу.

\begin{center}
  \renewcommand{\arraystretch}{1.5}
  \begin{tabularx}{\textwidth}{
    >{\centering\arraybackslash}p{0.15\textwidth}
    >{\centering\arraybackslash}p{0.25\textwidth}
    >{\centering\arraybackslash}p{0.25\textwidth}
    >{\centering\arraybackslash}X}
    \toprule
    \rowcolor{TableHeader}
    Ветвь & Кандидат & Условие ветви & Ограничение на \(a\) \\
    \midrule
    \(x\ge0\) & \answerblank[26mm] & \answerblank[25mm] & \answerblank[22mm] \\
    \addlinespace[0.25em]
    \(x<0\) & \answerblank[30mm] & \answerblank[25mm] & \answerblank[22mm] \\
    \bottomrule
  \end{tabularx}
\end{center}

\begin{checkpointbox}[Что происходит на границе \(a=-2\)?]
  Подставьте \(a=-2\) в обе формулы:
  \[
    x_1=\answerblank[15mm],
    \qquad
    x_2=\answerblank[15mm].
  \]
  Отрицательная ветвь требует \(x<0\). Выполняется ли это условие?
  \choicebox\ да \qquad \choicebox\ нет.

  Следовательно, при \(a=-2\) числитель имеет
  \answerblank[14mm] ноль, а два различных нуля возникают при
  \[
    a\ \answerblank[10mm]\ -2.
  \]
\end{checkpointbox}

\begin{reflectionbox}[Промежуточный вывод]
  Условие \(a>-2\) пока гарантирует только то, что
  \writinglines{2}
\end{reflectionbox}

\section{Проверяем найденные значения в знаменателе}

\begin{thinkbox}[Можно ли уже записывать ответ?]
  \choicebox\ Да. \qquad
  \choicebox\ Нет, потому что найденные значения могут
  \writinglines{1}
\end{thinkbox}

Обозначим знаменатель через
\[
  D(x)=x^2-2x-a.
\]

\stepheading{Шаг 1}{Проверяем \(x_1=a+2\)}

\begin{align*}
  D(x_1)
  &=(a+2)^2-2(a+2)-a\\
  &=\answerblank[42mm]\\
  &=\answerblank[32mm].
\end{align*}

Знаменатель равен нулю при
\[
  a=\answerblank[14mm]
  \qquad\text{или}\qquad
  a=\answerblank[14mm].
\]

\stepheading{Шаг 2}{Проверяем \(x_2=-\dfrac{a+2}{5}\)}

Начните подстановку и доведите выражение до произведения:
\[
  D(x_2)
  =\frac{(a+2)^2}{25}+\frac{2(a+2)}5-a.
\]
\workarea{33mm}

Получите:
\[
  D(x_2)=\frac{\answerblank[38mm]}{25}.
\]

Поэтому знаменатель равен нулю при
\[
  a=\answerblank[14mm]
  \qquad\text{или}\qquad
  a=\answerblank[14mm].
\]

\begin{fillbox}[Соберите проверку знаменателя]
  \begin{center}
    \renewcommand{\arraystretch}{1.45}
    \begin{tabularx}{\textwidth}{
      >{\centering\arraybackslash}p{0.24\textwidth}
      >{\centering\arraybackslash}p{0.38\textwidth}
      >{\centering\arraybackslash}X}
      \toprule
      \rowcolor{TableHeader}
      Кандидат & Знаменатель после подстановки & Исключаемые значения \(a\) \\
      \midrule
      \(x_1=a+2\) & \answerblank[38mm] & \answerblank[28mm] \\
      \addlinespace[0.25em]
      \(x_2=-\dfrac{a+2}{5}\) & \answerblank[43mm] & \answerblank[28mm] \\
      \bottomrule
    \end{tabularx}
  \end{center}
\end{fillbox}

\section{Собираем окончательный ответ}

Заполните две строки:
\[
  \text{числитель имеет два нуля при }a\in\answerblank[35mm];
\]
\[
  \text{нужно исключить }a\in\answerblank[47mm].
\]

На числовой прямой обведите исключаемые точки и выделите допустимые
промежутки.

\parameterstudentnumberlinefigure

\begin{resultbox}[Запишите ответ самостоятельно]
  \[
    a\in\answerblank[90mm].
  \]
\end{resultbox}

\begin{checkpointbox}[Проверка на трёх значениях]
  \begin{center}
    \renewcommand{\arraystretch}{1.45}
    \begin{tabularx}{\textwidth}{
      >{\centering\arraybackslash}p{0.12\textwidth}
      >{\centering\arraybackslash}p{0.28\textwidth}
      >{\centering\arraybackslash}p{0.36\textwidth}
      >{\centering\arraybackslash}X}
      \toprule
      \rowcolor{TableHeader}
      \(a\) & Нулей числителя & Что происходит со знаменателем & Корней дроби \\
      \midrule
      \(-2\) & \answerblank[12mm] & \answerblank[34mm] & \answerblank[12mm] \\
      \addlinespace[0.2em]
      \(-1\) & \answerblank[12mm] & \answerblank[34mm] & \answerblank[12mm] \\
      \addlinespace[0.2em]
      \(1\) & \answerblank[12mm] & \answerblank[34mm] & \answerblank[12mm] \\
      \bottomrule
    \end{tabularx}
  \end{center}
\end{checkpointbox}

\section{Почему условие \texorpdfstring{\(a>-2\)}{a > -2} видно на графике}

Положим \(b=a+2\). Тогда нули числителя определяются уравнением
\[
  3|x|-2x=b.
\]

\parametergraphworksheetfigure

\begin{fillbox}[Свяжите положение прямой и число пересечений]
  \begin{center}
    \renewcommand{\arraystretch}{1.45}
    \begin{tabularx}{0.92\textwidth}{@{}
      >{\centering\arraybackslash}p{0.16\textwidth}
      >{\centering\arraybackslash}p{0.24\textwidth}
      >{\centering\arraybackslash}p{0.2\textwidth}
      >{\centering\arraybackslash}X@{}}
      \toprule
      \rowcolor{TableHeader}
      Уровень & Условие на \(a\) & Пересечений & Нулей числителя \\
      \midrule
      \(b<0\) & \answerblank[18mm] & \answerblank[10mm] & \answerblank[10mm] \\
      \addlinespace[0.2em]
      \(b=0\) & \answerblank[18mm] & \answerblank[10mm] & \answerblank[10mm] \\
      \addlinespace[0.2em]
      \(b>0\) & \answerblank[18mm] & \answerblank[10mm] & \answerblank[10mm] \\
      \bottomrule
    \end{tabularx}
  \end{center}
\end{fillbox}

\begin{reflectionbox}[Ограничение графического рассуждения]
  Почему по этому графику нельзя увидеть исключаемые значения
  \(-1,0,3,8\)?
  \writinglines{2}
\end{reflectionbox}

\section{Более короткий способ: ищем общие корни}

Этот раздел показывает, как сразу найти параметры, при которых один из
нулей числителя запрещён знаменателем.

\begin{fillbox}[Составьте систему]
  Общий корень должен одновременно удовлетворять условиям
  \[
    \begin{cases}
      |3x|-2x-2-a=\answerblank[10mm],\\
      x^2-2x-a=\answerblank[10mm].
    \end{cases}
  \]
  Выразите \(a\) из каждого уравнения:
  \[
    a=\answerblank[42mm],
    \qquad
    a=\answerblank[38mm].
  \]
  После приравнивания получите уравнение без параметра:
  \[
    x^2=\answerblank[36mm].
  \]
\end{fillbox}

\begin{center}
  \renewcommand{\arraystretch}{1.5}
  \begin{tabularx}{\textwidth}{@{}
    >{\centering\arraybackslash}p{0.14\textwidth}
    >{\centering\arraybackslash}p{0.3\textwidth}
    >{\centering\arraybackslash}p{0.22\textwidth}
    >{\centering\arraybackslash}X@{}}
    \toprule
    \rowcolor{TableHeader}
    Случай & Уравнение & Общий корень \(x\) & Значение \(a\) \\
    \midrule
    \(x\ge0\) & \(x^2-3x+2=0\) & \answerblank[24mm] & \answerblank[24mm] \\
    \addlinespace[0.25em]
    \(x<0\) & \(x^2+3x+2=0\) & \answerblank[24mm] & \answerblank[24mm] \\
    \bottomrule
  \end{tabularx}
\end{center}

\begin{reflectionbox}[Смысл метода]
  Объясните, почему найденные значения параметра нужно не включить, а
  исключить из условия \(a>-2\).
  \writinglines{2}
\end{reflectionbox}

\section{Восстановите полное решение}

Закройте предыдущие страницы и запишите решение основной задачи по четырём
опорным пунктам:

\begin{practicebox}[Самостоятельная реконструкция]
  \begin{enumerate}[label=\textbf{\arabic*.}]
    \item Найдите два возможных нуля числителя и условия их ветвей.
    \item Определите, когда числитель имеет ровно два различных нуля.
    \item Проверьте оба значения в знаменателе.
    \item Объедините условия и запишите ответ.
  \end{enumerate}
  \writinglines{9}
\end{practicebox}

\section{Перенос метода}

\begin{practicebox}[Задание 1. Близкая модель]
  Найдите все значения \(a\), при которых уравнение
  \[
    \frac{|2x|-x-1-a}{x^2-x-a}=0
  \]
  имеет ровно два различных корня.
  \workarea{58mm}
  \longanswerline
\end{practicebox}

\begin{practicebox}[Задание 2. Анализ ошибки]
  Ученик получил условие \(a>-2\) и сразу записал его в ответ основной
  задачи.

  На каком этапе решение стало неполным? Является ли условие \(a>-2\)
  неверным или лишь недостаточным?
  \writinglines{3}
\end{practicebox}

\begin{reflectionbox}[Мой алгоритм]
  Закончите фразу: «Чтобы рациональное уравнение с параметром имело нужное
  количество корней, я сначала \ldots»
  \writinglines{4}
\end{reflectionbox}

\begin{questionsbox}[Перед завершением работы]
  \begin{itemize}[label=\(\square\)]
    \item Я сохранил условия ветвей при раскрытии модуля.
    \item Я различаю нули числителя и корни исходной дроби.
    \item Я проверил каждый кандидат в знаменателе.
    \item Я умею объяснить смысл условия \(a>-2\) графически.
    \item Я могу восстановить решение без готового образца.
  \end{itemize}
\end{questionsbox}
